\documentclass{main}

\title{Universal Inquiry}
\author{Gregory Lim}
\date{v1}

\begin{document}

\maketitle

\chapter{Logic}

A variable is an object.
If $a,b$ are variables, then $a\in b$ is a formula.
It is true if $a$ is in $b$, and false otherwise.
If $p$ is a formula, then $\neg p$ is a formula.
It is true if $p$ is false, and false otherwise.
If $p,q$ are formulas, then $p\land q$ is a formula.
It is true if $p,q$ are true, and false otherwise.
If $x$ is a variable, and $p$ is a formula, then $\exists x\,(p)$ is a formula.
It is true if there exists $x$ such that $p$ is true, and false otherwise.

\begin{notation}
$$
	p\lor q
	\equiv
	\neg(\neg p\land\neg q)
$$
\end{notation}

\begin{remark}
This says ``it is not the case that both $p$ and $q$ fail''.
\end{remark}

\begin{notation}
$$
	p\Rightarrow q
	\equiv
	\neg(p\land\neg q)
$$
\end{notation}

\begin{remark}
This says ``it is not the case that $p$ works but $q$ fails''.
\end{remark}

\begin{notation}
$$
	p\Leftrightarrow q
	\equiv
	(p\Rightarrow q)
	\land
	(q\Rightarrow p)
$$
\end{notation}

\begin{notation}
$$
	p\oplus q
	\equiv
	\neg(p\Leftrightarrow q)
$$
\end{notation}

\begin{remark}
This says ``it is not the case that $p$ and $q$ work the same''.
\end{remark}

\begin{notation}
$$
	\forall x\,(p)
	\equiv
	\neg\exists x\,(\neg p)
$$
\end{notation}

\begin{remark}
This says ``no counterexample exists''.
\end{remark}

\begin{notation}
$$
	X=Y
	\Leftrightarrow
	\forall a\,(
		a\in X
		\Leftrightarrow
		a\in Y
	)
	\land
	\forall B\,(
		X\in B
		\Leftrightarrow
		Y\in B
	)
$$
\end{notation}

\begin{remark}
This says ``$X$ and $Y$ look the same from the inside and outside''.
\end{remark}

\newpage

\begin{axiom}
$$
	\forall A,B\,(
		\forall x\,(
			x\in A
			\Leftrightarrow
			x\in B
		)
		\Rightarrow
		A=B
	)
$$
\end{axiom}

\begin{axiom}
$$
	\forall B,\exists U,\forall c\,(
		c\in U
		\Leftrightarrow
		\exists P\in B\,(
			c\in P
		)
	)
$$
\end{axiom}

\begin{definition}
$$
	U=\bigcup B
	\Leftrightarrow
	\forall c\,(
		c\in U
		\Leftrightarrow
		\exists P\in B\,(
			c\in P
		)
	)
$$
\end{definition}

\begin{definition}
$$
	A\cup B=\bigcup\{A,B\}
$$
\end{definition}

\begin{definition}
$$
	I=\bigcap B
	\Leftrightarrow
	\forall c\,(
		c\in I
		\Leftrightarrow
		\forall P\in B\,(
			c\in P
		)
	)
$$
\end{definition}

\begin{definition}
$$
	A\cap B=\bigcap\{A,B\}
$$
\end{definition}

\begin{axiom}
$$
	\forall X,
	\exists P,
	\forall S\,(
		S\in P
		\Leftrightarrow
		\forall s\in S\,(
			s\in X
		)
	)
$$
\end{axiom}

\begin{axiom}
$$
	\exists I\,(
		\varnothing\in I
		\land
		\forall i\in I\,(
			i\cup\{i\}\in I
		)
	)
$$
\end{axiom}

\begin{axiom}
$$
	\forall X\neq\varnothing,
	\exists x\in X\,(
		x\cap X
		=
		\varnothing
	)
$$
\end{axiom}

\begin{notation}
\begin{align*}
	\operatorname{free}(a\in b)
	&=
	\{a,b\}
	\\
	\operatorname{free}(\neg p)
	&=
	\operatorname{free}(p)
	\\
	\operatorname{free}(p\land q)
	&=
	\operatorname{free}(p)\cup\operatorname{free}(q)
	\\
	\operatorname{free}(\exists x\,(p))
	&=
	\operatorname{free}(p)\setminus\{x\}
\end{align*}
\end{notation}

\begin{axiom}
Let $B\not\in\operatorname{free}(p)$.
$$
\forall A\,(
	\forall x\in A,
	\exists!y\,(
		p
	)
	\Rightarrow
	\exists B, \forall y\,(
		y\in B
		\Leftrightarrow
		\exists x\in A\,(
			p
		)
	)
)
$$
\end{axiom}

\begin{definition}
$$
	A\times B
	=
	\{
		(a,b)
		\mid
		a\in A
		\land
		b\in B
	\}
$$
\end{definition}

\begin{definition}
$$
	Y^{X}
	=
	\{
		f\subseteq X\times Y
		\mid
		\forall x\in X,
		\exists!y\in Y\,(
			(x,y)\in f
		)
	\}
$$
\end{definition}

\end{document}
